\section{Appendix C: Implementation Code and User Interface}

\subsection{C.1 Software Architecture}

Our implementation follows modular software engineering principles:

\begin{itemize}
\item \textbf{Language:} Python 3.10
\item \textbf{Key Libraries:}
\begin{itemize}
\item \texttt{numpy}: Numerical computations
\item \texttt{pandas}: Data manipulation
\item \texttt{scipy}: Probability distributions, statistical tests
\item \texttt{networkx}: Graph algorithms
\item \texttt{docplex}: CPLEX Python API
\item \texttt{matplotlib}: Visualization
\item \texttt{geopandas}: GIS data handling
\end{itemize}
\item \textbf{Version Control:} Git (repository: github.com/Rolanduan/IMMC26601951)
\item \textbf{License:} MIT License (open source)
\end{itemize}

\subsection{C.2 Core Code Structure}

\subsubsection{Module Organization}

\begin{verbatim}
immc26601951/
├── main.py                 # Main execution script
├── config.py              # Configuration parameters
├── data/
│   ├── load_data.py       # Data loading and preprocessing
│   └── fit_distributions.py # Probability distribution fitting
├── models/
│   ├── deterministic_vrp.py    # Model 0
│   ├── stochastic_vrp.py       # Model 1 (chance-constrained)
│   ├── two_stage_vrp.py        # Model 2 (two-stage)
│   └── multi_depot_vrp.py      # Model 3
├── algorithms/
│   ├── alns.py            # Adaptive Large Neighborhood Search
│   ├── l_shaped.py        # Benders decomposition
│   └── clustering.py      # K-means clustering
├── evaluation/
│   ├── sensitivity.py     # Sensitivity analysis
│   └── validation.py      # Cross-validation
├── visualization/
│   ├── plot_routes.py     # GIS route visualization
│   └── plot_pareto.py     # Pareto frontier plotting
└── utils/
    ├── reliability.py     # Reliability calculation
    └── scenario_gen.py    # Scenario generation
\end{verbatim}

\subsection{C.3 Key Code Snippets}

\subsubsection{ALNS Main Loop}

\begin{verbatim}
class ALNS:
    def __init__(self, initial_solution, scenarios, alpha):
        self.solution = initial_solution
        self.scenarios = scenarios
        self.alpha = alpha
        self.destroy_operators = [self.random_removal,
                                  self.worst_removal,
                                  self.related_removal]
        self.repair_operators = [self.greedy_insertion,
                                 self.regret_insertion,
                                 self.reliability_insertion]
        self.operator_weights = self.initialize_weights()

    def solve(self, max_iterations=10000):
        best_solution = self.solution.copy()
        for iter in range(max_iterations):
            # Select operators
            destroy = self.select_operator(self.destroy_operators)
            repair = self.select_operator(self.repair_operators)

            # Generate new solution
            removed = destroy(self.solution, q=self.q)
            new_solution = repair(removed)

            # Evaluate
            if self.is_feasible(new_solution):
                if self.objective(new_solution) < self.objective(self.solution):
                    self.solution = new_solution
                    if self.objective(new_solution) < self.objective(best_solution):
                        best_solution = new_solution

            # Update weights
            self.update_weights(destroy, repair, success)

        return best_solution
\end{verbatim}

\subsubsection{Reliability Calculation}

\begin{verbatim}
def calculate_reliability(route, scenarios, T_max):
    """
    Calculate route reliability using scenario sampling.

    Parameters:
     - route: List of node indices [(i, j), ...]
    - scenarios: Dict of travel times {s: {(i,j): tau}}
    - T_max: Maximum allowed time

    Returns:
    - reliability: Probability of completing within T_max
    """
    on_time_count = 0
    for s in scenarios:
        total_time = sum(scenarios[s][edge] for edge in route)
        if total_time <= T_max:
            on_time_count += 1

    reliability = on_time_count / len(scenarios)
    return reliability
\end{verbatim}

\subsubsection{Scenario Generation}

\begin{verbatim}
import numpy as np
from scipy.stats import lognorm

def generate_scenarios(distributions, n_scenarios=100, seed=42):
    """
    Generate travel time scenarios from fitted distributions.

    Parameters:
    - distributions: Dict {(i,j): (shape, loc, scale) for lognormal}
    - n_scenarios: Number of scenarios to generate
    - seed: Random seed for reproducibility

    Returns:
    - scenarios: List of dicts {s: {(i,j): travel_time}}
    """
    np.random.seed(seed)
    scenarios = []

    for s in range(n_scenarios):
        scenario = {}
        for (i, j), params in distributions.items():
            shape, loc, scale = params
            # Sample from lognormal distribution
            travel_time = lognorm.rvs(shape, loc=loc, scale=scale)
            scenario[(i, j)] = travel_time
        scenarios.append(scenario)

    return scenarios
\end{verbatim}

\subsection{C.4 User Interface Design}

\subsubsection{Command-Line Interface}

\begin{verbatim}
# Generate routes for single-day delivery
python main.py --instance somalia_30 --model stochastic --alpha 0.85

# Run sensitivity analysis
python main.py --instance somalia_30 --sensitivity --parameter alpha

# Multi-objective optimization
python main.py --instance somalia_30 --multi-objective --epsilon-grid 50x50

# Visualize results
python visualization/plot_routes.py --solution routes_solution.pkl
\end{verbatim}

\subsubsection{Configuration File (config.yaml)}

\begin{verbatim}
# IMMC26601951 Configuration

problem:
  instance: "somalia_30"
  vehicles: 6
  capacity: 100
  time_windows: [8, 18]  # Operating hours

model:
  type: "stochastic"  # deterministic, stochastic, two_stage, multi_depot
  reliability_threshold: 0.85
  scenarios: 100
  validation_scenarios: 1000

algorithm:
  method: "alns"  # alns, cplex
  max_iterations: 10000
  convergence_gap: 0.01
  parallel_cores: 8

output:
  save_solution: true
  save_routes: true
  visualize: true
  export_format: ["pkl", "json", "csv"]
\end{verbatim}

\subsection{C.5 Data Input Format}

\subsubsection{Road Network CSV Format}

\begin{verbatim}
from_node,to_node,road_type,mean_travel_time,std_travel_time,elevation
0,1,paved,2.3,0.5,124
0,2,gravel,3.1,1.2,89
1,3,dirt,4.2,2.1,45
...
\end{verbatim}

\subsubsection{Demand CSV Format}

\begin{verbatim}
node_id,demand,time_window_start,time_window_end,priority
1,25,8,12,high
2,30,9,14,medium
3,15,10,16,low
...
\end{verbatim}

\subsection{C.6 Output Format}

\subsubsection{Solution JSON}

\begin{verbatim}
{
  "instance": "somalia_30",
  "model": "stochastic",
  "objective_value": 36.5,
  "reliability": 0.89,
  "routes": [
    {
      "vehicle_id": 1,
      "route": [0, 5, 12, 8, 3, 0],
      "load": 87,
      "expected_time": 6.2,
      "reliability": 0.91
    },
    ...
  ],
  "statistics": {
    "total_distance": 342.5,
    "total_time": 36.5,
    "fleet_utilization": 0.89,
    "unmet_demand": 0
  }
}
\end{verbatim}

\subsection{C.7 Web-Based User Interface (Conceptual Design)}

For non-technical users, we propose a web-based interface:

\subsubsection{Features:}
\begin{itemize}
\item \textbf{Interactive GIS Map:} Drag-and-drop warehouse locations, click to add customers
\item \textbf{Parameter Sliders:} Adjust $\alpha$, fleet size, capacity with live preview
\item \textbf{One-Click Optimization:} ``Generate Routes'' button
\item \textbf{Results Dashboard:} Key metrics (time, reliability, cost) with confidence intervals
\item \textbf{Scenario Analysis:} ``What-if'' tool for rainy season, security incidents
\end{itemize}

\subsubsection{Technology Stack:}
\begin{itemize}
\item Frontend: React + Leaflet (mapping)
\item Backend: FastAPI (Python web framework)
\item Database: PostgreSQL (PostGIS for spatial data)
\item Deployment: Docker containers (easy deployment)
\end{itemize}

\subsubsection{Screenshot (Conceptual):}
\begin{verbatim}
+----------------------------------------------------------+
|  Somalia Humanitarian Routing Optimizer                  |
+----------------------------------------------------------+
|  Map View              |  Parameters      |  Results    |
|                        |                  |             |
|  [Interactive GIS]     |  Alpha:  0.85    |  Time: 36.5h|
|  - Warehouse: 🏠      |  Fleet:    6     |  Rel:  89%  |
|  - Customer:  📍      |  Capacity: 100   |  Cost: $2120|
|  - Routes:    🛣️      |                  |             |
|                        | [Generate Routes]|            |
+----------------------------------------------------------+
\end{verbatim}

\subsection{C.8 Installation Instructions}

\begin{verbatim}
# Clone repository
git clone https://github.com/Rolanduan/IMMC26601951.git
cd IMMC26601951

# Create virtual environment
python -m venv venv
source venv/bin/activate  # On Windows: venv\Scripts\activate

# Install dependencies
pip install -r requirements.txt

# (Optional) Install CPLEX for exact methods
# Download from: https://www.ibm.com/products/ilog-cplex-optimization-studio

# Run example
python main.py --instance someria_30

# Run tests
pytest tests/
\end{verbatim}

\subsection{C.9 Dependencies}

\texttt{requirements.txt:}
\begin{verbatim}
numpy>=1.21.0
pandas>=1.3.0
scipy>=1.7.0
networkx>=2.6.0
matplotlib>=3.4.0
geopandas>=0.9.0
docplex>=2.0.0  # Optional: CPLEX interface
pytest>=6.2.0    # Testing
\end{verbatim}

\subsection{C.10 Performance Optimization Tips}

\subsubsection{Parallel Scenario Evaluation}

\begin{verbatim}
from multiprocessing import Pool

def evaluate_route_parallel(route, scenarios, T_max, n_cores=8):
    """Evaluate route reliability using parallel processing."""
    with Pool(n_cores) as pool:
        chunk_size = len(scenarios) // n_cores
        chunks = [scenarios[i:i+chunk_size]
                 for i in range(0, len(scenarios), chunk_size)]

        results = pool.starmap(calculate_reliability,
                              [(route, chunk, T_max) for chunk in chunks])

    return sum(results) / len(results)
\end{verbatim}

\subsubsection{Numba JIT Compilation}

\begin{verbatim}
from numba import jit

@jit(nopython=True)
def fast_objective_calculation(routes, travel_times, demands):
    """JIT-compiled objective function for speed."""
    total_time = 0.0
    for route in routes:
        for i in range(len(route) - 1):
            total_time += travel_times[route[i], route[i+1]]
    return total_time
\end{verbatim}

\subsection{C.11 Testing and Validation}

\subsubsection{Unit Tests}

\begin{verbatim}
# Test reliability calculation
def test_reliability_calculation():
    route = [(0, 1), (1, 2), (2, 0)]
    scenarios = [{(0,1): 2.0, (1,2): 3.0, (2,0): 2.5},
                 {(0,1): 2.5, (1,2): 3.5, (2,0): 3.0}]
    rel = calculate_reliability(route, scenarios, T_max=8.0)
    assert rel == 1.0  # Both scenarios under threshold

# Test scenario generation
def test_scenario_generation():
    dist = {(0,1): (0.5, 0, 2.0)}  # shape, loc, scale
    scenarios = generate_scenarios(dist, n_scenarios=100)
    assert len(scenarios) == 100
    assert all((0,1) in s for s in scenarios)
\end{verbatim}

\subsection{C.12 Future Development Roadmap}

\begin{enumerate}
\item \textbf{Version 2.0: Dynamic Re-optimization}
\begin{itemize}
\item Real-time data integration (weather, security)
\item Rolling horizon optimization
\item Mobile app for drivers
\end{itemize}

\item \textbf{Version 3.0: Machine Learning Enhancement}
\begin{itemize}
\item Learn travel time patterns from GPS data
\item Predictive modeling for disruptions
\item Adaptive reliability thresholds
\end{itemize}

\item \textbf{Version 4.0: Multi-Period Optimization}
\begin{itemize}
\item Inventory-routing problem (IRP)
\item Multi-day planning horizon
\item Warehouse inventory management
\end{itemize}
\end{enumerate}

\subsection{C.13 Contact and Support}

\begin{itemize}
\item \textbf{Repository:} github.com/Rolanduan/IMMC26601951
\item \textbf{Issues:} github.com/Rolanduan/IMMC26601951/issues
\item \textbf{Documentation:} github.com/Rolanduan/IMMC26601951/wiki
\item \textbf{License:} MIT License (free for academic and humanitarian use)
\end{itemize}

\textbf{Note:} Complete source code available in GitHub repository. This appendix provides key snippets for understanding implementation approach and reproducing results.
